\chapter{Resumen}

Este Trabajo Fin de Grado desarrolla el banco de ensayos de un mecanismo de elevación basado en una
\emph{caracola}: un tambor de radio variable que, accionado por un servomotor paso a paso de lazo
cerrado, enrolla un cable y levanta una barra articulada con una masa de \SI{2}{\kilo\gram}. La idea
que se quiere contrastar es la del \emph{fusee} de la relojería: si el radio del tambor crece donde
la tensión del cable es pequeña y decrece donde es grande, el par que ve el motor se mantiene
aproximadamente constante a lo largo del recorrido.

Para poder contrastarla se ha construido una plataforma completa: una placa adaptadora diseñada en
KiCad alrededor de una Raspberry Pi Pico~2, con bus RS-485, entradas digitales protegidas, entradas
analógicas acondicionadas, salidas de potencia y dos zócalos mikroBUS; un firmware en C++ que
gobierna el accionamiento por Modbus RTU, lee la célula de carga a través de un acondicionador
ZSC31014, hace el \emph{homing} con dos finales de carrera y ejecuta ciclos automáticos
B$\rightarrow$A$\rightarrow$B; un protocolo binario propio con CRC-8; y una aplicación web (FastAPI
y WebSocket) con tres páginas: monitor en tiempo real, configuración del mecanismo con simulador, y
visor de ensayos que superpone la curva teórica a cada medida.

El banco muestrea a \SI{20.0}{\hertz} con una dispersión del periodo de \SI{0.06}{\milli\second},
graba cada ensayo en cuentas brutas y repite el recorrido con una desviación típica de
\SI{11}{mrev} de motor. Se han analizado dos campañas de ensayos, la segunda con el
ciclo de ensayo reorganizado para que arranque en el extremo inferior del recorrido.

El resultado principal es metodológico y, en parte, negativo. El modelo estático reproduce bien la
forma de la señal de la célula de carga ($R^2 = 0{,}97$), pero se demuestra que esa comparación es
insuficiente: al no estar calibrada la célula, el ajuste absorbe un cambio de signo y el mismo
$R^2$ admite dos geometrías opuestas. El par del accionamiento, que no necesita calibración, deshace
la ambigüedad, y al hacerlo contradice la configuración vigente del mecanismo: sobre el tramo del
recorrido en que el cable trabaja con carga, el modelo predice que la tensión \emph{baja} un
\SI{20}{\percent} mientras que el par gravitatorio medido \emph{sube} un \SI{56}{\percent}. El radio
efectivo deducido crece un \SI{93}{\percent}, cuando para mantener el par constante bastaría con que
creciera un \SI{24}{\percent}. La caracola, tal y como está descrita en la configuración, no
compensa: sobrecompensa, y con el sentido invertido respecto del que pide el modelo.

El trabajo entrega, por tanto, un banco de ensayos reproducible, un modelo verificado en su
formulación y un procedimiento de contraste, junto con un diagnóstico claro de lo que falta para
cerrar la validación: medir las cotas reales del mecanismo y calibrar la célula con masas patrón.

\paragraph{Palabras clave:} tambor de radio variable, caracola, compensación de par, servomotor paso
a paso, Modbus RTU, RS-485, célula de carga, ZSC31014, Raspberry Pi Pico 2, adquisición de datos.

\chapter{Abstract}

This Bachelor's Thesis develops the test bench for a lifting mechanism based on a variable-radius
spiral drum (a \emph{fusee}-like drum, here called \emph{caracola}): driven by a closed-loop stepper
servomotor, it winds a cable that lifts a hinged bar carrying a \SI{2}{\kilo\gram} mass. The idea
under test is the one behind the horological fusee: if the drum radius grows where the cable tension
is small and shrinks where it is large, the motor torque stays roughly constant along the stroke.

A complete platform was built to test it: a KiCad-designed adapter board around a Raspberry Pi
Pico~2, with an RS-485 bus, protected digital inputs, conditioned analogue inputs, high-side outputs
and two mikroBUS sockets; C++ firmware that commands the drive over Modbus RTU, reads the load cell
through a ZSC31014 conditioner, performs homing with two limit switches and runs automatic
B$\rightarrow$A$\rightarrow$B cycles; a custom binary protocol with CRC-8; and a web application
(FastAPI and WebSocket) with three pages: real-time monitor, mechanism configuration with a
simulator, and a test viewer that overlays the theoretical curve on each measurement.

The bench samples at \SI{20.0}{\hertz} with a period spread of \SI{0.06}{\milli\second}, logs every
test in raw counts and repeats the stroke with a standard deviation of \SI{11}{mrev} of
motor shaft. Two test campaigns were analysed, the second one with the test cycle reorganised to
start at the lower end of the stroke.

The main result is methodological and partly negative. The static model reproduces the shape of the
load-cell signal well ($R^2 = 0.97$), but that comparison is shown to be insufficient: since the
cell is uncalibrated, the fit absorbs a sign change and the same $R^2$ admits two opposite
geometries. The drive torque, which needs no calibration, resolves the ambiguity --- and in doing so
contradicts the current configuration of the mechanism: over the loaded part of the stroke the model
predicts the tension to \emph{fall} by \SI{20}{\percent} while the measured gravitational torque
\emph{rises} by \SI{56}{\percent}. The implied effective radius grows by \SI{93}{\percent}, whereas
a \SI{24}{\percent} growth would suffice for constant torque. The spiral drum, as configured, does
not compensate: it overcompensates, and in the opposite sense to the one the model requires.

The thesis therefore delivers a reproducible test bench, a model verified in its formulation and a
contrast procedure, together with a clear diagnosis of what remains in order to close the
validation: measuring the real dimensions of the mechanism and calibrating the load cell with
reference masses.

\paragraph{Keywords:} variable radius drum, fusee, torque compensation, closed-loop stepper servo,
Modbus RTU, RS-485, load cell, ZSC31014, Raspberry Pi Pico 2, data acquisition.
